% !PRINTED 86
\weberpage{86}
%<<CONT

two rows coincide, while if \(k>\nu\), it vanishes by the hypothesis, since (12) is then a \((\nu+1)\)-rowed determinant of the matrix (3). Thus from (11), since \(A\) is different from zero, we obtain
\begin{equation}
        \sum_{\mu=1}^{n}a_\mu^{(k)}x_\mu=0,\qquad k=1,2,\ldots,m,
\tag{13}
\end{equation}
that is, the system of equations (1) is satisfied by (7).

We shall single out from the theorem just proved the special case in which \(m=n-1\) and \(\nu=n-1\). In this case only \gesperrt{one} of the unknowns remains arbitrary, and the ratios of the unknowns are completely determined. We may state this result as follows:

Denote the \((n-1)\)-rowed determinants of the matrix
\begin{equation}
\begin{array}{cccc}
 a_1^{(1)}&a_2^{(1)}&\cdots&a_n^{(1)}\\
 a_1^{(2)}&a_2^{(2)}&\cdots&a_n^{(2)}\\
 \vdots&\vdots&&\vdots\\
 a_1^{(n-1)}&a_2^{(n-1)}&\cdots&a_n^{(n-1)}
\end{array}
\tag{14}
\end{equation}
taken with alternating sign by \(A_1,A_2,\ldots,A_n\), and assume that at least one of these quantities is different from zero. Then the solution of the system
\begin{equation}
\begin{aligned}
 a_1^{(1)}x_1+a_2^{(1)}x_2+\cdots+a_n^{(1)}x_n&=0,\\
 a_1^{(2)}x_1+a_2^{(2)}x_2+\cdots+a_n^{(2)}x_n&=0,\\
 &\vdots\\
 a_1^{(n-1)}x_1+a_2^{(n-1)}x_2+\cdots+a_n^{(n-1)}x_n&=0
\end{aligned}
\tag{15}
\end{equation}
is given by the ratios
\begin{equation}
        x_1:x_2:\cdots:x_n=A_1:A_2:\cdots:A_n.
\tag{16}
\end{equation}
Thus for \(n=3\) we obtain the solution of the system of equations, often occurring in geometry,
\begin{equation}
        ax+by+cz=0,\qquad a'x+b'y+c'z=0
\tag{17}
\end{equation}
in the form
\begin{equation}
        x:y:z=bc'-cb':ca'-ac':ab'-ba'.
\tag{18}
\end{equation}
